% =====================================================================
%  CAPÍTULO 2 — MARCO DEL TRABAJO
% =====================================================================
\chapter{Marco del trabajo}
\label{cap:marco}

Antes de exponer cómo se ha analizado la comunicación de sostenibilidad de las empresas
europeas conviene delimitar el terreno sobre el que se asienta el trabajo. Ese terreno tiene
dos dimensiones que se entrelazan. La primera es normativa: la información de sostenibilidad
no se publica en el vacío, sino bajo un marco regulatorio europeo que en los últimos años ha
cambiado de manera radical y que condiciona qué deben contar las empresas y con qué nivel de
exigencia. La segunda es teórica: para interpretar por qué las empresas comunican como
comunican —y por qué a veces el discurso se aleja de la práctica— resulta imprescindible
apoyarse en la literatura sobre legitimidad, gestión de impresiones y \textit{greenwashing}.
A ambas dimensiones se dedica este capítulo, que se cierra con una panorámica del estado de la
cuestión en la aplicación de la inteligencia artificial al análisis de los informes
corporativos.

\section{Marco normativo: de la NFRD a la CSRD}
\label{sec:normativa}

La regulación europea de la información no financiera nació, en su forma moderna, con la
Directiva 2014/95/UE, comúnmente conocida como NFRD por sus siglas inglesas
(\textit{Non-Financial Reporting Directive}). Esta norma obligó por primera vez a las grandes
empresas europeas de más de quinientos empleados a incorporar en su informe de gestión una
serie de contenidos relativos al medio ambiente, a las cuestiones sociales y laborales, al
respeto de los derechos humanos, a la lucha contra la corrupción y a la diversidad en sus
órganos de administración \parencite{nfrd2014}. Fue, sin duda, un avance: situó la
sostenibilidad en el terreno de la rendición de cuentas obligatoria. Sin embargo, su diseño
adolecía de debilidades que el tiempo hizo evidentes. Al no imponer un formato común, los
informes resultaban enormemente heterogéneos y difícilmente comparables entre sí; no exigía
verificación externa de la información, de modo que su fiabilidad descansaba en la buena
voluntad del emisor; y, sobre todo, no incorporaba el principio de doble materialidad ni
descendía al detalle de indicadores concretos. El resultado era un reporting con frecuencia
genérico, voluntarista y orientado a la imagen más que a la rendición de cuentas. A ello se
sumaba un alcance limitado, que cubría a unas once mil empresas en el conjunto de la Unión.

La respuesta a esas limitaciones llegó con la Directiva (UE) 2022/2464, la
\textit{Corporate Sustainability Reporting Directive} (CSRD), que no se limita a retocar la
NFRD sino que la refunde y amplía sustancialmente \parencite{csrd2022}. Tres son sus
innovaciones de mayor calado. La primera es la consagración del principio de
\textbf{doble materialidad}, en virtud del cual la empresa debe informar tanto de los impactos
que su actividad genera sobre la sociedad y el medio ambiente —materialidad de impacto— como
de los riesgos y oportunidades que las cuestiones de sostenibilidad representan para su propio
negocio —materialidad financiera—. La segunda es la estandarización obligatoria de los
contenidos a través de los \textit{European Sustainability Reporting Standards} (ESRS), que
sustituyen la libertad de formato de la NFRD por un lenguaje común y comparable. La tercera es
la exigencia de aseguramiento externo, inicialmente con un nivel de garantía limitada, que
somete la información de sostenibilidad a un escrutinio independiente análogo —aunque todavía
no equivalente— al de la información financiera. A todo ello se añade la obligación de un
etiquetado digital de los datos para facilitar su tratamiento automatizado, y una ampliación
del alcance que multiplica por más de cuatro el número de empresas afectadas, hasta rondar las
cincuenta mil en toda la Unión.

El despliegue de la CSRD es progresivo. El primer grupo en aplicarla —las grandes empresas que
ya estaban sujetas a la NFRD— ha tenido que reportar conforme al nuevo marco a partir del
ejercicio 2024, con publicación de los informes a lo largo de 2025; los siguientes escalones
incorporan, en años sucesivos, a otras grandes empresas no cubiertas anteriormente, a las
pequeñas y medianas cotizadas y, finalmente, a determinadas empresas de fuera de la Unión con
actividad significativa en ella. Esta gradualidad es precisamente lo que confiere valor
analítico a la ventana temporal estudiada en este trabajo: los ejercicios 2022 y 2023
pertenecen todavía, mayoritariamente, al régimen de la NFRD, mientras que 2024 constituye el
primer ejercicio plenamente sujeto a la CSRD para las grandes empresas de interés público que
componen la muestra. El contraste entre ambos regímenes, sintetizado en la
tabla~\ref{tab:nfrd-csrd}, es el eje sobre el que se articula la cuarta pregunta de
investigación.

\begin{table}[H]
\centering
\caption{Principales diferencias entre la NFRD y la CSRD}
\label{tab:nfrd-csrd}
\begin{tabular}{lll}
\toprule
\textbf{Dimensión} & \textbf{NFRD (2014/95/UE)} & \textbf{CSRD (2022/2464/UE)} \\
\midrule
Alcance & $\sim$11.000 empresas & $\sim$50.000 empresas \\
Estándares & Sin formato único & ESRS obligatorios \\
Doble materialidad & No & Sí \\
Aseguramiento externo & No requerido & Garantía limitada obligatoria \\
Formato digital & No & XBRL/iXBRL obligatorio \\
Información de la cadena de valor & Mínima & Extensa \\
Objetivos climáticos & No específicos & Alineados con el Acuerdo de París \\
\bottomrule
\end{tabular}
\\[4pt]
{\footnotesize Fuente: elaboración propia a partir de \textcite{nfrd2014} y \textcite{csrd2022}.}
\end{table}

El corazón técnico de la CSRD reside en los ESRS, aprobados mediante el Reglamento Delegado
(UE) 2023/2772 \parencite{esrs2023}. Se trata de doce estándares que organizan toda la
información de sostenibilidad. Dos de ellos son transversales y de aplicación obligatoria en
todos los casos: el ESRS 1 fija los requisitos generales y la arquitectura del sistema
—el propio concepto de doble materialidad, los horizontes temporales, la debida diligencia—,
mientras que el ESRS 2 establece las divulgaciones generales sobre gobernanza, estrategia,
gestión de riesgos e impactos, y métricas y objetivos. Los diez estándares restantes son
temáticos y se aplican únicamente en la medida en que la empresa los considere materiales.
Cinco son ambientales y reproducen, no por casualidad, los objetivos de la Taxonomía de la
Unión: el cambio climático (ESRS E1), la contaminación (E2), los recursos hídricos y marinos
(E3), la biodiversidad y los ecosistemas (E4) y el uso de recursos y la economía circular
(E5). Cuatro son sociales y recorren el círculo concéntrico de personas afectadas por la
actividad: la plantilla propia (S1), los trabajadores de la cadena de valor (S2), las
comunidades afectadas (S3) y los consumidores y usuarios finales (S4). El último, de
gobernanza, atañe a la conducta empresarial (G1): anticorrupción, cultura ética, prácticas de
pago y protección de los denunciantes. Esta taxonomía de doce estándares no es un mero detalle
regulatorio: como se verá, constituye la rejilla con la que se ha medido la cobertura temática
de los informes analizados.

El marco europeo se completa con dos piezas que, aun sin ser el objeto directo del análisis,
condicionan el comportamiento de las empresas. La primera es la Taxonomía de la Unión, el
Reglamento (UE) 2020/852, que establece un sistema de clasificación de las actividades
económicas ambientalmente sostenibles en torno a seis objetivos —mitigación y adaptación al
cambio climático, agua, economía circular, contaminación y biodiversidad— y sujeto a tres
criterios: contribuir sustancialmente a alguno de ellos, no perjudicar significativamente a
los demás (el principio \textit{Do No Significant Harm}) y respetar unas salvaguardas sociales
mínimas \parencite{taxonomia2020}. Las grandes empresas publican desde hace algunos ejercicios
el porcentaje de su cifra de negocio, de sus inversiones y de sus gastos operativos que está
alineado con la Taxonomía, un dato que se ha convertido en un indicador de referencia. La
segunda pieza es el Reglamento (UE) 2019/2088, conocido como SFDR, que impone obligaciones de
transparencia a los participantes en los mercados financieros —gestoras, aseguradoras, fondos
de pensiones— y que, aunque no se dirige directamente a las empresas no financieras, ejerce
sobre ellas una presión indirecta considerable: los inversores institucionales necesitan datos
de sostenibilidad de calidad sobre las compañías de su cartera para poder cumplir con sus
propias obligaciones de divulgación \parencite{sfdr2019}. El resultado es un ecosistema
regulatorio en el que la demanda de información de sostenibilidad fiable y comparable proviene
simultáneamente del regulador y del mercado.

\section{Marco teórico: comunicación corporativa, legitimidad y \textit{greenwashing}}
\label{sec:teoria}

La existencia de un marco normativo cada vez más exigente no garantiza, por sí sola, que la
comunicación de las empresas sea veraz o sustancial. Para entender por qué el discurso
corporativo puede divergir de la práctica conviene recurrir a la teoría de la legitimidad,
que constituye el sustrato conceptual de buena parte de la investigación contable sobre
sostenibilidad. Según esta perspectiva, las organizaciones operan bajo un contrato implícito
con la sociedad y necesitan que su actividad sea percibida como legítima, esto es, conforme a
los valores y expectativas del entorno. Cuando ese contrato se ve amenazado —por una crisis,
por un escándalo ambiental o, simplemente, por una mayor exigencia social—, la empresa
despliega estrategias para restaurar o reforzar su legitimidad. La comunicación de
sostenibilidad es una de las herramientas más visibles de ese esfuerzo, y de ahí su ambivalencia:
puede reflejar un compromiso real o puede operar como un instrumento de gestión de la imagen.

\textcite{cho2015} llevaron esta intuición un paso más allá al introducir el concepto de
\textit{fachada organizativa} (\textit{organizational façade}). En su análisis, los informes
de sostenibilidad funcionan a menudo como una construcción discursiva orientada a satisfacer
las presiones institucionales sin que medie un cambio efectivo en las prácticas de la empresa.
Los autores hablan de una \textit{hipocresía organizada} en la que el discurso y la acción se
desacoplan de forma deliberada. La aportación es fundamentalmente conceptual —no ofrece
métricas— pero proporciona el marco interpretativo para una de las hipótesis centrales de este
trabajo: la posibilidad de que un lenguaje más optimista enmascare un desempeño menos
sustancial.

Si Cho y sus colegas describen el fenómeno en términos generales, \textcite{hahn2014}
descienden al detalle de los mecanismos concretos. A partir de un análisis cualitativo de
informes elaborados conforme al estándar GRI \parencite{gri2021}, identifican seis estrategias mediante las que
las empresas legitiman los aspectos negativos de su desempeño: la abstracción —declaraciones
vagas que evitan cuantificar—, el desplazamiento temporal —remitir la solución a un futuro
indeterminado—, la externalización de responsabilidades, la comparación favorable, la negación
del carácter negativo del impacto y la invocación de acciones compensatorias. Estas estrategias
son especialmente relevantes para esta investigación porque dos de ellas, la abstracción y el
desplazamiento temporal, admiten una operacionalización cuantitativa directa: la primera se
traduce en la ausencia de cifras concretas, y la segunda, en un lenguaje prospectivo de
promesas sin respaldo numérico. Como se detallará en el capítulo de metodología, ambas se han
convertido en componentes del índice de \textit{greenwashing} construido en este trabajo.

La pregunta sobre la calidad de la divulgación había sido abordada empíricamente por
\textcite{michelon2015}, cuyo análisis arroja una conclusión incómoda: las prácticas que
habitualmente se asocian a un reporting de mayor calidad —la publicación de informes
independientes, el uso del estándar GRI o la verificación externa— no garantizan, por sí
mismas, una información más sustancial. En muchos casos, sostienen los autores, el reporting es
más simbólico que real. Su contribución más valiosa para este trabajo es, sin embargo,
metodológica: descomponen la calidad de la divulgación en dimensiones medibles, entre las que
destacan la amplitud —número de temas tratados—, la profundidad —nivel de detalle—, la
especificidad —el peso de lo cuantitativo frente a lo cualitativo— y la neutralidad —la
inclusión, o no, de información desfavorable—. La especificidad, en particular, es la dimensión
que vertebra la parte central del análisis empírico de esta investigación.

El instrumental para medir el lenguaje sobre el texto financiero procede de un trabajo ya
clásico. \textcite{loughran2011} demostraron que los diccionarios de sentimiento de uso
general clasifican erróneamente numerosas palabras cuando se aplican al ámbito financiero:
términos como \textit{liability} (pasivo), \textit{tax} (impuesto) o \textit{cost} (coste),
negativos en un diccionario psicológico convencional, son neutros en el lenguaje de los
negocios. Para resolverlo construyeron un diccionario específico para textos financieros que
clasifica las palabras en categorías como negativo, positivo, incertidumbre, lenguaje
litigioso o modalidad —fuerte y débil—. Estas categorías permiten medir no solo el tono, sino
también el grado de vaguedad o cautela de un texto: las palabras de incertidumbre y de
modalidad débil son, en la práctica, indicadores del lenguaje evasivo (\textit{hedging}) que
Hahn y Lülfs describían como estrategia de abstracción. El diccionario de Loughran y McDonald
es, por ello, una de las herramientas empleadas en este trabajo, si bien con una cautela que
se discute más adelante: fue concebido para los informes anuales estadounidenses y su
traslación a informes de sostenibilidad europeos exige prudencia.

El punto de encuentro entre toda esta tradición teórica y las técnicas de inteligencia
artificial lo marca el trabajo de \textcite{bingler2022}, que constituye el referente
metodológico más directo de esta investigación. Aplicando un modelo de lenguaje especializado
en texto climático —ClimateBERT— a más de mil quinientos informes alineados con las
recomendaciones del TCFD \parencite{tcfd2017}, los autores muestran que buena parte de las divulgaciones climáticas
corporativas constituyen \textit{cheap talk}: compromisos retóricos que no se acompañan de
acciones concretas. Detectan, además, un patrón de \textit{cherry-picking}, consistente en
divulgar de forma destacada la información climática que no resulta comprometida —la que no
cuesta nada— mientras se omite la material. Su hallazgo más relevante para el presente trabajo
es que las divulgaciones voluntarias se asocian a un mayor grado de \textit{cheap talk} que las
obligatorias, lo que sugiere de forma natural una hipótesis sobre el tránsito de la NFRD a la
CSRD. La presente investigación adopta su lógica —medir señales de discurso vacío directamente
sobre el texto— y la extiende a un marco más amplio: no solo el clima, sino el conjunto de las
dimensiones ESRS, y no solo un régimen, sino la transición entre dos.

\section{Estado de la cuestión: el PLN aplicado al reporting de sostenibilidad}
\label{sec:estado}

La aplicación de técnicas de procesamiento del lenguaje natural al análisis de los informes
corporativos ha conocido en los últimos años una expansión notable, impulsada por la
disponibilidad de modelos de lenguaje cada vez más capaces y por la creciente presión
regulatoria sobre la información de sostenibilidad. El trabajo de Bingler y sus colegas, ya
comentado, representa el ejemplo más influyente de esta corriente, pero no el único. La
investigación se ha movido en dos grandes direcciones complementarias. La primera es el
análisis de sentimiento y de especificidad, que busca cuantificar el tono y la concreción del
discurso, y que se apoya en modelos especializados como FinBERT, para el lenguaje financiero,
o las distintas variantes de ClimateBERT, para el climático. La segunda es el modelado de
temas, que pretende descubrir, sin imponerlos de antemano, los asuntos de los que hablan los
informes, y que ha evolucionado desde los métodos probabilísticos clásicos hacia enfoques
basados en representaciones semánticas del texto.

Una línea más reciente, y especialmente pertinente para este trabajo, es la que propone medir
la cobertura de los estándares ESRS mediante diccionarios. \textcite{suta2025} desarrollan un
enfoque basado en diccionarios para evaluar en qué medida los informes cubren los distintos
temas de divulgación de los ESRS, y validan así una estrategia que combina la transparencia y
la reproducibilidad de los métodos léxicos con la pertinencia normativa del marco europeo.
Esta aportación respalda directamente una de las decisiones metodológicas de la presente
investigación, que construye un diccionario propio de términos ESRS para medir la amplitud
temática de cada informe.

Conviene, no obstante, leer esta literatura con sentido crítico. Los trabajos pioneros, como
el de Bingler y sus colegas, se circunscriben al ámbito climático y a un periodo anterior a la
plena vigencia de la CSRD; los análisis de calidad de la divulgación, como el de Michelon y
sus colegas, son anteriores a la irrupción de las técnicas de PLN y descansan en la
codificación manual; y los enfoques basados en diccionarios, aun siendo transparentes,
capturan la presencia de un vocabulario, no su profundidad ni su veracidad. Cada método ilumina
una faceta del problema y deja otras en sombra. De esa constatación nace la estrategia central
de este trabajo, que no se confía a una única técnica sino que triangula varias —diccionarios,
modelado de temas y modelos de lenguaje— para obtener una imagen más robusta.

El hueco que esta investigación pretende cubrir se sitúa, así, en la intersección de varias
carencias de la literatura previa. Ninguno de los trabajos revisados analiza el conjunto del
STOXX Europe 600 —es decir, la gran empresa cotizada europea en su diversidad sectorial y
geográfica— bajo el marco completo de la CSRD y los ESRS, combinando de forma sistemática el
modelado de temas, el análisis de sentimiento y especificidad y la detección de señales de
\textit{greenwashing}, y haciéndolo además sobre una ventana temporal que captura el tránsito
de un régimen regulatorio voluntario a uno obligatorio. La tabla~\ref{tab:matriz} resume la
posición de los principales antecedentes y el lugar que ocupa este trabajo respecto de ellos.

\begin{table}[H]
\centering
\caption{Matriz de literatura y aportación del presente trabajo}
\label{tab:matriz}
\begin{tabular}{p{3.1cm}p{2.4cm}p{3.0cm}p{4.0cm}}
\toprule
\textbf{Autor (año)} & \textbf{Muestra} & \textbf{Método} & \textbf{Gap que deja} \\
\midrule
Loughran \& McDonald (2011) & Informes 10-K (EE.\,UU.) & Diccionario, análisis textual & No sostenibilidad, no europeo \\
Hahn \& Lülfs (2014) & DJIA + DAX & Análisis cualitativo & No cuantitativo, muestra pequeña \\
Michelon et al. (2015) & Empresas cotizadas & Análisis de contenido manual & Anterior a la CSRD, sin PLN \\
Cho et al. (2015) & Empresas cotizadas & Análisis conceptual & Sin operacionalización \\
Bingler et al. (2022) & $>$1.500 informes TCFD & ClimateBERT & Solo clima, anterior a la CSRD \\
Suta et al. (2025) & Informes ESRS & Diccionarios ESRS & Enfoque léxico, sin triangulación \\
\textbf{Este TFG (2025)} & \textbf{STOXX 600} & \textbf{Temas + sentimiento + \gwindex} & \textbf{Cubre el marco CSRD/ESRS completo} \\
\bottomrule
\end{tabular}
\\[4pt]
{\footnotesize Fuente: elaboración propia a partir de la revisión de la literatura.}
\end{table}
